\documentclass[12pt,a4paper,onecolumn]{exam}
\usepackage{amsmath}
\usepackage{amssymb}
\usepackage{graphicx}
\usepackage{float}
\usepackage{booktabs}
\usepackage{geometry}
\usepackage{tikz}
\usepackage[skins]{tcolorbox} % Use [skins] to get the rounded shape


% --- Define our simple \questionheader command ---
\newcommand{\questionheader}[1]{%
  \begin{tcolorbox}[
    enhanced,
    colback=black,
    coltext=white,
    boxrule=0pt,
    fontupper=\Large\bfseries,
    arc=4mm
  ]
  #1
  \end{tcolorbox}%
}
% --- End of definition ---
\newtcolorbox{answerbox}[1]{
  boxrule=0.4pt,   % Sets the thickness of the border
  colback=white,   % Sets the background color
  height=#1,       % <-- Use the first argument (#1) as the height
}


\usepackage[font = small]{caption}
\usepackage{subcaption}

\newenvironment{blanksolution}
  {%
    \renewcommand{\solutiontitle}{\noindent}%
    \begin{solution}%
  }%
  {\end{solution}}
\usepackage{listings}

\lstset{
    language=Python,
    basicstyle=\ttfamily,
    }

\pagestyle{headandfoot}
% \firstpageheader{Assignment 2}{}{Dwaipayan Haldar}
\runningheader{Assignment 8}{Signal Processing in Practise}{Dwaipayan Haldar}

\begin{document}

\begingroup
    \centering
    \LARGE E9 222 Signal Processing in Practice\\
    \LARGE Assignment 8\\[0.5em]
    \large \today\\[0.5em]
    \large Dwaipayan Haldar \par
\endgroup
\noindent\rule{\textwidth}{0.5pt}
\printanswers
\renewcommand{\solutiontitle}{\noindent\textbf{Ans:}\enspace}
\setlength{\parskip}{1pt}


%% ─────────────────────────────────────────────────────────────────────────────
\questionheader{Problem 1: MixUp — Interpolating Data for Robust Augmentation}

\begin{solution}

\noindent\textbf{Architecture:} SimpleCNN — Conv1(3$\to$32) $\to$ BN $\to$ ReLU $\to$ Pool $\to$ Conv2(32$\to$64) $\to$ BN $\to$ ReLU $\to$ Pool $\to$ FC1(4096$\to$128) $\to$ BN $\to$ ReLU $\to$ FC2(128$\to$10). Trained with Adam (lr=0.001), Cross-Entropy loss, 20 epochs on CIFAR-10 (normalised to $[-1,1]$).

\vspace{4pt}
\noindent\textbf{Accuracy Summary:}

\begin{center}
\resizebox{\linewidth}{!}{%
\begin{tabular}{lc}
\toprule
\textbf{Method} & \textbf{Test Accuracy} \\
\midrule
1.1 Raw (no augmentation) & 0.7392 \\
1.2 All augmentations (HFlip + Crop + ColorJitter + Gaussian Noise) & 0.7214 \\
\midrule
1.3 Mixup $\alpha=0.1$ & 0.7334 \\
1.3 Mixup $\alpha=0.2$ & 0.7346 \\
1.3 Mixup $\alpha=0.4$ & \textbf{0.7556} \\
1.3 Mixup $\alpha=1.0$ & 0.7530 \\
\bottomrule
\end{tabular}}
\end{center}

\vspace{4pt}
\noindent\textbf{Problem 1.1 — Raw Training:}
\begin{figure}[H]
    \centering
    \includegraphics[width=0.55\linewidth]{loss_curves_raw.png}
    \caption{Training and validation loss — raw (no augmentation).}
\end{figure}

\noindent\textbf{Problem 1.2 — All Augmentations:}
\begin{figure}[H]
    \centering
    \includegraphics[width=0.55\linewidth]{loss_curves_all_aug.png}
    \caption{Training and validation loss — all augmentations combined.}
\end{figure}

\noindent\textbf{Problem 1.3 — Mixup ($\lambda \sim \mathrm{Beta}(\alpha,\alpha)$):}

\begin{figure}[H]
\centering
\begin{subfigure}[b]{0.48\linewidth}
    \includegraphics[width=\linewidth]{loss_curves_mixup_alpha_0.1.png}
    \caption{$\alpha=0.1$}
\end{subfigure}
\begin{subfigure}[b]{0.48\linewidth}
    \includegraphics[width=\linewidth]{loss_curves_mixup_alpha_0.2.png}
    \caption{$\alpha=0.2$}
\end{subfigure}\\[4pt]
\begin{subfigure}[b]{0.48\linewidth}
    \includegraphics[width=\linewidth]{loss_curves_mixup_alpha_0.4.png}
    \caption{$\alpha=0.4$}
\end{subfigure}
\begin{subfigure}[b]{0.48\linewidth}
    \includegraphics[width=\linewidth]{loss_curves_mixup_alpha_1.0.png}
    \caption{$\alpha=1.0$}
\end{subfigure}
\caption{Mixup training and validation loss for different $\alpha$ values. Higher $\alpha$ raises train loss (stronger mixing) but narrows the train/val gap.}
\end{figure}

\begin{figure}[H]
    \centering
    \includegraphics[width=\linewidth]{mixup_visualization.png}
    \caption{Mixup augmentation visualisation for $\alpha \in \{0.1, 0.2, 0.4, 1.0\}$. Each row shows the same 8 images blended with a random partner; $\alpha=1.0$ produces the most aggressive blending.}
\end{figure}

\noindent\textbf{Observations:}
\begin{itemize}\setlength\itemsep{1pt}
    \item The raw data clearly overfits. The validation curve diverges while the training loss converges. The augmentation reduces overfitting and give a more generalized performance. 
    \item Traditional augmentation (1.2) \emph{hurts} accuracy vs.\ the raw baseline — aggressive spatial/colour transforms at training time increase noise without providing the interpolation-based regularisation that Mixup offers.
    \item Mixup consistently \emph{outperforms} the raw baseline. Best accuracy is at $\alpha=0.4$ (0.7556), a $\sim$1.6\,pp gain over raw.
    \item Small $\alpha$ ($0.1$, $0.2$) keeps $\lambda$ near 0 or 1, so mixing is mild — train loss stays low but val loss is higher (less regularisation). Larger $\alpha$ ($0.4$, $1.0$) increases the train/val loss gap but the \emph{val loss is lower}, indicating better generalisation.
    \item $\alpha=1.0$ (Uniform mixing) performs slightly worse than $\alpha=0.4$, suggesting over-regularisation for this shallow network.
\end{itemize}

\end{solution}


%% ─────────────────────────────────────────────────────────────────────────────
\questionheader{Problem 2: Manifold Mixup — Interpolation in Feature Space}

\begin{solution}

\noindent Manifold Mixup applies $\tilde{g}_k = \lambda\, g_k(x_i) + (1-\lambda)\, g_k(x_j)$, $\tilde{y} = \lambda y_i + (1-\lambda)y_j$ where $g_k$ is the output of the $k$-th hidden layer. The same SimpleCNN is used with $\alpha=0.2$; the final classification layer (FC2) is excluded from mixing.

\vspace{4pt}
\noindent\textbf{Layer key:} $k=1$: after Conv1+Pool; \;$k=2$: after Conv2+Pool; \;$k=3$: after FC1; \;$k=\text{random}$: uniformly random $k\in\{1,2,3\}$ per batch.

\begin{center}
\begin{tabular}{lc}
\toprule
\textbf{Mixing Layer} & \textbf{Test Accuracy} \\
\midrule
$k=1$ (after Conv1) & \textbf{0.7522} \\
$k=2$ (after Conv2) & 0.7419 \\
$k=3$ (after FC1)   & 0.7429 \\
$k=\text{random}$   & 0.7399 \\
\bottomrule
\end{tabular}
\end{center}

\begin{figure}[H]
\centering
\begin{subfigure}[b]{0.48\linewidth}
    \includegraphics[width=\linewidth]{loss_curves_manifold_k1.png}
    \caption{$k=1$ (after Conv1+Pool)}
\end{subfigure}
\begin{subfigure}[b]{0.48\linewidth}
    \includegraphics[width=\linewidth]{loss_curves_manifold_k2.png}
    \caption{$k=2$ (after Conv2+Pool)}
\end{subfigure}\\[4pt]
\begin{subfigure}[b]{0.48\linewidth}
    \includegraphics[width=\linewidth]{loss_curves_manifold_k3.png}
    \caption{$k=3$ (after FC1)}
\end{subfigure}
\begin{subfigure}[b]{0.48\linewidth}
    \includegraphics[width=\linewidth]{loss_curves_manifold_krandom.png}
    \caption{$k=\text{random}$}
\end{subfigure}
\caption{Training and validation loss for Manifold Mixup at each layer setting ($\alpha=0.2$).}
\end{figure}

\noindent\textbf{Observations:}
\begin{itemize}\setlength\itemsep{1pt}
    \item Mixing at the earliest hidden layer ($k=1$) gives the best accuracy (0.7522), comparable to input-space Mixup ($\alpha=0.4$: 0.7556). Earlier mixing acts closer to input Mixup, preserving low-level feature diversity.
    \item Mixing at deeper layers ($k=2$, $k=3$) reduces accuracy slightly — at those stages, feature maps are more task-specific, so interpolation in that space is less natural and may produce out-of-manifold representations.
    \item Random $k$ per batch performs worst among the four, likely because inconsistent mixing depth prevents the model from adapting its intermediate representations stably.
    \item All Manifold Mixup variants outperform traditional augmentation (0.7214) and are competitive with raw training (0.7392).
\end{itemize}

\end{solution}

\end{document}
